\chapter{Product types}

In this chapter product types will be discussed. A product type is called as
such because it can be seen as a Cartesian product of other types.

There are three product types:
\begin{itemize}
	\item \textbf{Tuple} (\see{sec:tuples})
	\item \textbf{Record} (\see{sec:records})
	\item \textbf{Struct} (\see{sec:structs})
\end{itemize}

All product types follow the same rules of the default (\see{ch:defaults}),
hence, when a variable or any other value of such type is declared, it is
initialized with its zero value.\footnotemark[1]

\footnotetext[1]{There is one exception: see~\See{sec:recursive_structs}.}

\section{Tuples}
\label{sec:tuples}

A Tuple is an ordered collection of zero or more heterogeneous values.

\begin{minted}{\giovilexer}
someTuple tuple<Text, Bool, Real>
someTuple <- <"foo", true, 420.69>
\end{minted}

Tuples being ordered means that two tuple types with the same fields (same
count and equally typed), but in different position, are considered different.

\section{Records}
\label{sec:records}

A Record is an unordered collection of zero or more named (labeled)
heterogeneous values.

\begin{minted}{\giovilexer}
someRecord record<field1 Text, field2 Bool, field3 Real>
someRecord <- <field1:"foo", field2:true, field3:420.69>
\end{minted}

Records being unordered means that two record types with the same fields (same
count, equally named and equally typed) are equal, regardless of the order
their fields appear at.

\section{Structs}
\label{sec:structs}

A Struct is a named unordered collection of zero or more named (labeled)
heterogeneous fields. In other words, a struct type is a type you can define
and later refer to via its name. Structs could be seen as Records
(\see{sec:records}) with a typename associated to it.

Before being able to use a struct type, it must first be declared.

\begin{minted}{\giovilexer}
struct SomeStruct <field1 Text, field2 Bool, field3 Real>
someStruct SomeStruct
\end{minted}

\paragraph{Struct constants} Like for functions (see
\See{sec:struct_functions}), a struct type can also contain constant
declarations, that can later be accessed via the dot operator.

\subsection{Recursive struct types}
\label{sec:recursive_structs}

Recursive struct type definitions are allowed by using the \code{self} keyword
in place of the field type.

\todo{menzionare RFC di larval types e uso del typename}

\begin{minted}{\giovilexer}
struct Node <
	data Text,
	nextNode self
>
\end{minted}

All types in Giovi have a zero value (see~\See{ch:defaults}) and all values
of such type are initialized with it, and the same applies to structs;
however recursive structs are the exception.\\*
In the case of the a recursive struct, the fields of such type will be
initialized with a lazy approach: the initialization of such fields happens
only upon the first access/use.

\begin{minted}{\giovilexer}
struct Node <
	data Text default "default_data",
	nextNode self
>
root Node <- <data:"root_data">
root -> _out       # <"root_data", _>
root.nextNode -> _ # nextNode is now initialized
root -> _out       # <"root_data", <"default_data", _>>
\end{minted}

\clearpage

\subsection{Aggregation and composition}

When modelling a certain domain, it is often needed to define an aggregation or
composition relationship between types. To achieve aggregation, practically
speaking, it's simply possible to define a field of a previously typed.

\begin{minted}{\giovilexer}
struct Fullname <name Text, surname Text>
struct Person <age Int, fullname Fullname>
Person -> _out # Person<age Int, fullname Fullname>
\end{minted}

At times, however, it may be desirable to "flatten/inline" such aggregated
fields; instead of manually copying the fields in the new declaration, it's
possible to specify the typename of the other type you want to aggregate.

\begin{minted}{\giovilexer}
struct Person <age Int> + Fullname
Person -> _out # <age Int, name Text, surname Text>
\end{minted}

The aggregate type can also be specified before the new declaration; by doing
so, the fields will appear before the ones in the new declaration.

\begin{minted}{\giovilexer}
Fullname + struct Person <age Int>
Person -> _out # <name Text, surname Text, age Int>
\end{minted}

Multiple aggregations are possible

\begin{minted}{\giovilexer}
Vehicle, Engine + struct Car<brand Text> + AirCondition, ReclineableSeats
Car -> _out
# Car<
#	Vehicle fields...,
#	Engine fields...,
#	brand Text,
#	AirCondition fields...,
#	ReclineableSeats fields...,
#>
\end{minted}

\todo{Formalizzare aggregazioni multiple, overrides e defaults}

